\chapter{Magnesium blood test}

\section*{Definition and Overview}
The magnesium blood test is a clinical diagnostic test that measures the concentration of magnesium in a patient's serum or plasma. Magnesium is the fourth most abundant cation in the human body and the second most abundant intracellular cation. It plays a vital role as a cofactor in over 300 enzymatic reactions, including those regulating protein synthesis, muscle and nerve function, blood glucose control, and blood pressure. While a standard serum magnesium test is the most common method to evaluate magnesium status, it represents only about 1\% of the body's total magnesium stores, as most magnesium is stored in bone and intracellular compartments, meaning clinical correlation is essential to diagnose subtle deficiencies or excesses.

\section*{Epidemiology}
Hypomagnesaemia (low magnesium levels) is a common electrolyte imbalance, particularly among hospitalised patients. It affects approximately 2\% to 15\% of the general population, but the prevalence rises to 50\% to 60\% in intensive care unit (ICU) patients and individuals with chronic alcoholism. Hypermagnesaemia (high magnesium levels) is much rarer, primarily occurring in patients with acute or chronic kidney disease who ingest magnesium-containing medications, such as antacids or laxatives. In Australia, electrolyte panel monitoring is a routine component of pathology testing in acute and primary care settings, with magnesium tests commonly ordered for patients presenting with unexplained neuromuscular or cardiovascular symptoms.

\section*{Aetiology and Pathophysiology}
Fluctuations in serum magnesium levels reflect a disruption in the balance between gastrointestinal absorption, renal excretion, and intracellular-extracellular shifts.
\begin{itemize}
    \item \textbf{Gastrointestinal malabsorption:} Chronic diarrhoea, malabsorptive syndromes (such as coeliac disease or Crohn's disease), or poor dietary intake reducing magnesium absorption in the small intestine.
    \item \textbf{Renal wasting:} Drug-induced renal loss (e.g., from loop diuretics, thiazides, proton pump inhibitors, or aminoglycosides) or genetic tubular disorders (like Gitelman or Bartter syndromes).
    \item \textbf{Endocrine and metabolic influences:} Uncontrolled diabetes mellitus (where osmotic diuresis increases renal magnesium excretion) or hyperaldosteronism.
    \item \textbf{Intracellular shifts:} Rapid uptake of magnesium into cells, which can occur during the treatment of diabetic ketoacidosis or in hungry bone syndrome following parathyroidectomy.
    \item \textbf{Renal insufficiency:} Decreased glomerular filtration rate (GFR) in advanced kidney disease impairs the kidneys' ability to excrete magnesium, leading to systemic accumulation.
    \item \textbf{Exogenous overload:} Excessive administration of magnesium sulfate (e.g., in pre-eclampsia treatment) or excessive ingestion of magnesium-containing antacids or laxatives.
\end{itemize}

\section*{Clinical Features}
The clinical manifestations of abnormal magnesium levels range from asymptomatic mild fluctuations to severe, life-threatening neurological and cardiovascular crises.
\begin{itemize}
    \item \textbf{Hypomagnesaemia symptoms (mild to moderate):} Muscle cramps, weakness, tremors, hyperreflexia, fasciculations, and paresthesias.
    \item \textbf{Hypomagnesaemia signs (severe):} Tetany, positive Chvostek and Trousseau signs, generalised seizures, and cardiac arrhythmias, including Torsades de Pointes.
    \item \textbf{Hypermagnesaemia signs (mild to moderate):} Nausea, vomiting, flushing, headache, lethargy, and diminished deep tendon reflexes.
    \item \textbf{Hypermagnesaemia signs (severe):} Somnolence, muscle paralysis, severe hypotension, bradycardia, complete heart block, and cardiac arrest.
    \item \textbf{Associated electrolyte imbalances:} Concurrent hypocalcaemia (due to impaired parathyroid hormone secretion and action) and hypokalaemia, which are refractory to correction until magnesium is replenished.
\end{itemize}

\section*{Differential Diagnosis}
Because the symptoms of magnesium imbalances are non-specific, they must be differentiated from other metabolic and neurological conditions.
\begin{itemize}
    \item \textbf{Hypocalcaemia or hypercalcaemia:} Electrolyte disorders that present with similar neuromuscular irritability (tetany, spasms) or lethargy, and frequently coexist with magnesium anomalies.
    \item \textbf{Hypokalaemia or hyperkalaemia:} Potassium disturbances that present with muscle weakness, ECG changes, and cramping.
    \item \textbf{Primary neurological disorders:} Epilepsy or peripheral neuropathies that present with seizures or sensory disturbances in the absence of electrolyte abnormalities.
    \item \textbf{Myasthenia gravis:} A neuromuscular disease causing weakness and fatigue that can be exacerbated by magnesium administration, which inhibits acetylcholine release.
    \item \textbf{Chronic fatigue syndrome:} Presents with persistent fatigue and muscle aches but with normal serum magnesium and electrolyte levels.
\end{itemize}

\section*{When to Seek Medical Care}
Individuals with signs of severe muscle irritability, weakness, or unexplained heart palpitations should undergo medical evaluation, including electrolyte testing.
\textbf{Call 000 (or local emergency services) for:}
\begin{itemize}
    \item Sudden onset of muscle paralysis or extreme, generalised muscle weakness.
    \item Seizures or loss of consciousness.
    \item Chest pain, severe palpitations, or irregular heartbeat.
    \item Difficulty breathing or severe lethargy in a patient receiving intravenous magnesium therapy.
\end{itemize}
Other reasons to seek medical care include persistent muscle twitching, chronic cramps, unexplained numbness, or if you are taking long-term proton pump inhibitors and experience fatigue or muscle spasms.

\section*{Diagnosis and Investigations}
The diagnosis of magnesium disorders requires quantitative laboratory measurement of serum magnesium levels, supplemented by additional biochemical and cardiac tests.
\begin{itemize}
    \item \textbf{Serum magnesium test:} The standard diagnostic test; the normal reference range is typically 0.70 to 1.10 mmol/L, though this varies slightly by laboratory.
    \item \textbf{Fractional excretion of magnesium (FEMg):} Measured using concurrent spot urine and blood samples to differentiate renal wasting from gastrointestinal loss.
    \item \textbf{Serum potassium and calcium levels:} Co-ordered to identify concurrent hypokalaemia or hypocalcaemia, which frequently complicate hypomagnesaemia.
    \item \textbf{Electrocardiogram (ECG):} Essential to assess for electrical conduction abnormalities, such as prolonged QT or PR intervals, widening of the QRS complex, or T-wave changes.
    \item \textbf{Renal function tests (eGFR, Urea, Creatinine):} Evaluated to check for renal impairment, which is the primary risk factor for hypermagnesaemia.
\end{itemize}

\section*{Management}
The management of magnesium imbalances focuses on correcting the serum level safely, treating the underlying cause, and avoiding toxicity during replacement.
\begin{itemize}
    \item \textbf{Oral magnesium replacement:} Used for mild, asymptomatic hypomagnesaemia, utilising organic salts (e.g., magnesium aspartate or citrate) which are better absorbed than magnesium oxide, as the oxide can cause diarrhoea.
    \item \textbf{Intravenous magnesium sulfate:} Reserved for severe or symptomatic hypomagnesaemia, particularly if seizures, tetany, or arrhythmias are present, administered slowly with close monitoring of reflexes and vitals.
    \item \textbf{Addressing the underlying cause:} Reviewing and discontinuing offending drugs, managing gastrointestinal disease, or controlling diabetes.
    \item \textbf{Hypermagnesaemia management (mild):} Stopping all exogenous magnesium sources and administering intravenous fluids with loop diuretics to promote renal clearance.
    \item \textbf{Intravenous calcium gluconate:} Administered in severe hypermagnesaemia as an immediate antagonist to magnesium's neuromuscular effects, protecting the heart while renal excretion or haemodialysis is initiated.
\end{itemize}

\section*{Complications}
\begin{itemize}
    \item \textbf{Refractory hypokalaemia and hypocalcaemia:} Failure to correct potassium and calcium levels due to unaddressed hypomagnesaemia, prolonging neuromuscular symptoms.
    \item \textbf{Cardiac arrhythmias:} Development of lethal arrhythmias, including ventricular tachycardia, ventricular fibrillation, or Torsades de Pointes.
    \item \textbf{Neuromuscular block and respiratory depression:} Occurs in severe hypermagnesaemia, leading to diaphragmatic paralysis and respiratory failure.
    \item \textbf{Hypotension and bradycardia:} Severe cardiovascular instability associated with magnesium toxicity.
    \item \textbf{Diarrhoea and gastrointestinal distress:} Common side effects of excessive oral magnesium supplementation, which can worsen dehydration and electrolyte depletion.
\end{itemize}

\section*{Prognosis and Outcomes}
The prognosis for patients with abnormal magnesium levels is generally excellent when the condition is recognised early and corrected appropriately. Mild hypomagnesaemia responds rapidly to oral therapy, although long-term maintenance may be required if the underlying cause (such as malabsorption or medication use) cannot be resolved. Hypermagnesaemia has a favourable outcome if renal function is intact and exogenous intake is stopped. However, unrecognised severe electrolyte imbalances carry significant risks of cardiac arrest and neurological complications in acute care settings.

\section*{Prevention and Self-Care}
\begin{itemize}
    \item Consume a diet rich in magnesium-containing foods, including leafy green vegetables, nuts, seeds, whole grains, and legumes.
    \item Avoid overusing magnesium-containing over-the-counter medications, such as certain antacids or laxatives, especially if you have known kidney issues.
    \item Monitor your symptoms and have regular blood tests if you are prescribed long-term loop diuretics or proton pump inhibitors.
    \item Limit alcohol consumption, as chronic heavy drinking impairs renal conservation of magnesium and leads to depletion.
    \item Take magnesium supplements only under the guidance of a healthcare professional to avoid gastrointestinal side effects or toxicity.
\end{itemize}

\section*{Sources and Further Reading}
The following reputable organisations provide further clinical details and guidelines regarding magnesium and electrolyte blood tests:
\begin{itemize}
    \item Pathology Tests Explained (Australia). \textit{Detailed guide to the magnesium blood test.} https://pathologytestsexplained.org.au/
    \item Healthdirect Australia. \textit{Role of magnesium, deficiency symptoms, and testing.} https://www.healthdirect.gov.au/
    \item Mayo Clinic. \textit{Understanding electrolyte panels and magnesium status.} https://www.mayoclinic.org/
    \item MSD Manual (Professional Version). \textit{Pathophysiology and management of hypomagnesemia and hypermagnesemia.} https://www.msdmanuals.com/
    \item Cleveland Clinic. \textit{Overview of magnesium levels and health effects.} https://my.clevelandclinic.org/
\end{itemize}
